% コマ15: 継承の連鎖を final/mycc.py の 1 ファイルへ畳み込む
% 各回の mycc.py は前回クラスの差分でしかない
\documentclass[border=8pt]{standalone}
% 講義資料の図(materials/sessions/figures/*.tex)共通プリアンブル
% 各図の .tex から % 講義資料の図(materials/sessions/figures/*.tex)共通プリアンブル
% 各図の .tex から % 講義資料の図(materials/sessions/figures/*.tex)共通プリアンブル
% 各図の .tex から \input{figure-preamble} で読み込む。
% 配色・フォントは latex/session-template.tex と揃える。

% 日本語(LuaLaTeX + luatexja)。等幅セル内の日本語も和文フォントで出す
\usepackage{luatexja}
\usepackage{luatexja-fontspec}
\setmainfont{Noto Sans CJK JP}
\setmainjfont{Noto Sans CJK JP}
\setmonofont{Inconsolata}[Scale=MatchLowercase]
\setmonojfont{Noto Sans CJK JP}

\usepackage{xcolor}
\definecolor{TcAccentDark}{HTML}{583B66}
\definecolor{TcAccentLight}{HTML}{EEE6F2}
\definecolor{TcAccentLighter}{HTML}{F7F3F9}
\definecolor{TcEmph}{HTML}{E64B17}
\definecolor{TcText}{HTML}{262626}
\definecolor{TcGrayText}{HTML}{737373}
\definecolor{TcGrayBg}{HTML}{F2F2F2}

\usepackage{tikz}
\usetikzlibrary{arrows.meta,positioning,calc,decorations.pathreplacing,fit,backgrounds}

\tikzset{
  % テキスト
  every node/.style={text=TcText},
  annot/.style={font=\small, text=TcText, align=left},
  gannot/.style={font=\small, text=TcGrayText, align=center},
  addr/.style={font=\ttfamily\small, text=TcGrayText},
  % メモリセル
  memcell/.style={draw=TcAccentDark, line width=0.6pt, fill=white,
                  minimum height=8.5mm, minimum width=40mm,
                  font=\ttfamily, inner sep=2pt, outer sep=0pt},
  memcell gray/.style={memcell, fill=TcGrayBg},
  memcell hi/.style={memcell, fill=TcAccentLighter},
  memcell dashed/.style={memcell, dashed, fill=none, text=TcGrayText},
  % 領域(セクション図など)
  region/.style={draw=TcAccentDark, line width=0.6pt, fill=white,
                 minimum width=48mm, align=center, inner sep=3pt, outer sep=0pt},
  % 制御フロー図のボックス
  cfgbox/.style={draw=TcAccentDark, line width=0.6pt, fill=white, rounded corners=1.5mm,
                 minimum width=28mm, minimum height=8mm, font=\small, align=center,
                 inner sep=2pt},
  % 矢印
  ptr/.style={-{Stealth[length=2.8mm]}, line width=1.1pt, color=TcEmph},
  flow/.style={-{Stealth[length=2.4mm]}, line width=0.8pt, color=TcAccentDark},
  flow back/.style={flow, dashed},
  regptr/.style={-{Stealth[length=2.8mm]}, line width=1.1pt, color=TcAccentDark},
  % 補助
  bracestyle/.style={decorate, decoration={brace, amplitude=4pt}, line width=0.8pt,
                     color=TcAccentDark},
}
 で読み込む。
% 配色・フォントは latex/session-template.tex と揃える。

% 日本語(LuaLaTeX + luatexja)。等幅セル内の日本語も和文フォントで出す
\usepackage{luatexja}
\usepackage{luatexja-fontspec}
\setmainfont{Noto Sans CJK JP}
\setmainjfont{Noto Sans CJK JP}
\setmonofont{Inconsolata}[Scale=MatchLowercase]
\setmonojfont{Noto Sans CJK JP}

\usepackage{xcolor}
\definecolor{TcAccentDark}{HTML}{583B66}
\definecolor{TcAccentLight}{HTML}{EEE6F2}
\definecolor{TcAccentLighter}{HTML}{F7F3F9}
\definecolor{TcEmph}{HTML}{E64B17}
\definecolor{TcText}{HTML}{262626}
\definecolor{TcGrayText}{HTML}{737373}
\definecolor{TcGrayBg}{HTML}{F2F2F2}

\usepackage{tikz}
\usetikzlibrary{arrows.meta,positioning,calc,decorations.pathreplacing,fit,backgrounds}

\tikzset{
  % テキスト
  every node/.style={text=TcText},
  annot/.style={font=\small, text=TcText, align=left},
  gannot/.style={font=\small, text=TcGrayText, align=center},
  addr/.style={font=\ttfamily\small, text=TcGrayText},
  % メモリセル
  memcell/.style={draw=TcAccentDark, line width=0.6pt, fill=white,
                  minimum height=8.5mm, minimum width=40mm,
                  font=\ttfamily, inner sep=2pt, outer sep=0pt},
  memcell gray/.style={memcell, fill=TcGrayBg},
  memcell hi/.style={memcell, fill=TcAccentLighter},
  memcell dashed/.style={memcell, dashed, fill=none, text=TcGrayText},
  % 領域(セクション図など)
  region/.style={draw=TcAccentDark, line width=0.6pt, fill=white,
                 minimum width=48mm, align=center, inner sep=3pt, outer sep=0pt},
  % 制御フロー図のボックス
  cfgbox/.style={draw=TcAccentDark, line width=0.6pt, fill=white, rounded corners=1.5mm,
                 minimum width=28mm, minimum height=8mm, font=\small, align=center,
                 inner sep=2pt},
  % 矢印
  ptr/.style={-{Stealth[length=2.8mm]}, line width=1.1pt, color=TcEmph},
  flow/.style={-{Stealth[length=2.4mm]}, line width=0.8pt, color=TcAccentDark},
  flow back/.style={flow, dashed},
  regptr/.style={-{Stealth[length=2.8mm]}, line width=1.1pt, color=TcAccentDark},
  % 補助
  bracestyle/.style={decorate, decoration={brace, amplitude=4pt}, line width=0.8pt,
                     color=TcAccentDark},
}
 で読み込む。
% 配色・フォントは latex/session-template.tex と揃える。

% 日本語(LuaLaTeX + luatexja)。等幅セル内の日本語も和文フォントで出す
\usepackage{luatexja}
\usepackage{luatexja-fontspec}
\setmainfont{Noto Sans CJK JP}
\setmainjfont{Noto Sans CJK JP}
\setmonofont{Inconsolata}[Scale=MatchLowercase]
\setmonojfont{Noto Sans CJK JP}

\usepackage{xcolor}
\definecolor{TcAccentDark}{HTML}{583B66}
\definecolor{TcAccentLight}{HTML}{EEE6F2}
\definecolor{TcAccentLighter}{HTML}{F7F3F9}
\definecolor{TcEmph}{HTML}{E64B17}
\definecolor{TcText}{HTML}{262626}
\definecolor{TcGrayText}{HTML}{737373}
\definecolor{TcGrayBg}{HTML}{F2F2F2}

\usepackage{tikz}
\usetikzlibrary{arrows.meta,positioning,calc,decorations.pathreplacing,fit,backgrounds}

\tikzset{
  % テキスト
  every node/.style={text=TcText},
  annot/.style={font=\small, text=TcText, align=left},
  gannot/.style={font=\small, text=TcGrayText, align=center},
  addr/.style={font=\ttfamily\small, text=TcGrayText},
  % メモリセル
  memcell/.style={draw=TcAccentDark, line width=0.6pt, fill=white,
                  minimum height=8.5mm, minimum width=40mm,
                  font=\ttfamily, inner sep=2pt, outer sep=0pt},
  memcell gray/.style={memcell, fill=TcGrayBg},
  memcell hi/.style={memcell, fill=TcAccentLighter},
  memcell dashed/.style={memcell, dashed, fill=none, text=TcGrayText},
  % 領域(セクション図など)
  region/.style={draw=TcAccentDark, line width=0.6pt, fill=white,
                 minimum width=48mm, align=center, inner sep=3pt, outer sep=0pt},
  % 制御フロー図のボックス
  cfgbox/.style={draw=TcAccentDark, line width=0.6pt, fill=white, rounded corners=1.5mm,
                 minimum width=28mm, minimum height=8mm, font=\small, align=center,
                 inner sep=2pt},
  % 矢印
  ptr/.style={-{Stealth[length=2.8mm]}, line width=1.1pt, color=TcEmph},
  flow/.style={-{Stealth[length=2.4mm]}, line width=0.8pt, color=TcAccentDark},
  flow back/.style={flow, dashed},
  regptr/.style={-{Stealth[length=2.8mm]}, line width=1.1pt, color=TcAccentDark},
  % 補助
  bracestyle/.style={decorate, decoration={brace, amplitude=4pt}, line width=0.8pt,
                     color=TcAccentDark},
}

\begin{document}
\begin{tikzpicture}[node distance=0mm,
    stg/.style={cfgbox, minimum width=84mm, minimum height=7.5mm, font=\small,
                align=left, text width=78mm},
    cls/.style={font=\ttfamily\small},
    note/.style={annot, text=TcGrayText, align=left}]

  \node[gannot, anchor=south] at (38mm, 6mm) {通常回（コマ2〜コマ14）};

  \foreach \i/\y/\t in {
      0/0/{{\ttfamily Codegen02}  算術式・関数プロローグ／エピローグ},
      1/-11/{{\ttfamily Codegen03}  {\ttfamily \_locals}・{\ttfamily codegen\_lval\_Var}・代入},
      2/-22/{{\ttfamily Codegen04}  ラベル生成・{\ttfamily if}／{\ttfamily else}・比較演算},
      3/-33/{{\ttfamily Codegen05}  {\ttfamily while}／{\ttfamily for}・{\ttfamily break}／{\ttfamily continue}},
      4/-44/{{\ttfamily Codegen06}  関数定義・関数呼び出し・再帰},
      5/-55/{{\ttfamily Codegen07}  {\ttfamily codegen\_lval} の分離・{\ttfamily \&}／{\ttfamily *}},
      6/-66/{{\ttfamily Codegen08}  {\ttfamily ty\_str}・{\ttfamily \_load\_ty}／{\ttfamily \_store\_ty}},
      7/-77/{{\ttfamily Codegen09}  ポインタ演算・添字・{\ttfamily sizeof}・{\ttfamily malloc}},
      8/-88/{{\ttfamily Codegen10}  {\ttfamily .data}・文字列ラベル・{\ttfamily printf}},
      9/-99/{{\ttfamily Codegen11}  式の走査の網羅},
      10/-110/{{\ttfamily Codegen12}  {\ttfamily struct}・{\ttfamily .}／{\ttfamily ->}・レイアウト},
      11/-121/{{\ttfamily Codegen13}  {\ttfamily \_globals}・{\ttfamily .bss}・{\ttfamily \&\&}／{\ttfamily ||}／{\ttfamily !}},
      12/-132/{{\ttfamily Codegen14}  {\ttfamily \#define}・複数ファイル}} {
    \node[stg] (b\i) at (38mm, \y mm) {\t};
  }

  \foreach \i [evaluate=\i as \j using int(\i+1)] in {0,...,11} {
    \draw[flow] (b\i.south) -- (b\j.north);
  }

  \node[note, anchor=east] at (-6mm, -66mm)
      {importlib 継承\\（各回は差分だけを書く）};

  % ---- 全体を束ねて final へ ----
  \draw[bracestyle] ($(b0.north east)+(4mm,0)$) -- ($(b12.south east)+(4mm,0)$);

  \node[cfgbox, minimum width=52mm, minimum height=44mm, align=center, font=\small]
      (final) at (140mm, -66mm)
      {{\ttfamily final/mycc.py}\\[2mm]
       継承なし。\\
       全ハンドラを 1 クラスに展開する};
  \draw[ptr] (88mm, -66mm) -- (final.west);

  % ---- 検証フロー ----
  \node[cfgbox, minimum width=34mm, minimum height=8mm, font=\footnotesize]
      (tests) at (140mm, -102mm) {final/tests/（fixed17）};
  \node[cfgbox, minimum width=34mm, minimum height=8mm, font=\footnotesize]
      (runner) at (140mm, -118mm) {test\_runner.py};
  \draw[flow] (final.south) -- (tests.north);
  \draw[flow] (tests.south) -- (runner.north);
  \node[gannot, anchor=west, align=left] at (160mm, -118mm)
      {17 本すべてが通れば\\標準トラック完成の目安};

  % ---- 補足 ----
  \node[gannot, anchor=north, align=center] at (76mm, -146mm)
      {各回の {\ttfamily mycc.py} は前の回のクラスを継承した差分でしかないので、
       そのまま {\ttfamily final/} へ置いても動かない。\\
       連鎖をたどって、どの回の何がどこに効いているかを確かめながら 1 つのクラスへまとめる
       （{\ttfamily tokenize} / {\ttfamily parse} も同じ連鎖で引き継がれている）};
\end{tikzpicture}
\end{document}
